% Auto-generated from raw.txt
% Usage:
%   \vocabentry{word}{/ipa/ (pos)}{definition}{example}{note}

\vocabentry{Globalization}
{/ˌɡloʊbələˈzeɪʃn/ (n)}
{The process by which businesses or other organizations develop international influence or start operating on an international scale.}
{Globalization has led to a significant increase in international trade and cultural exchange.}
{Đây là một danh từ dùng để chỉ quá trình toàn cầu hóa, sự kết nối đa quốc gia về kinh tế, văn hóa và chính trị.}

\vocabentry{Integration}
{/ˌɪntɪˈɡreɪʃn/ (n)}
{The action or process of integrating or combining several things.}
{The economic integration of European nations has created a powerful unified market.}
{Đây là một danh từ dùng để chỉ sự hội nhập, đặc biệt là hội nhập kinh tế và thị trường toàn cầu.}

\vocabentry{Multinational}
{/ˌmʌltiˈnæʃnəl/ (adj)}
{Including or involving several countries or individuals of several nationalities.}
{Multinational corporations play a central role in the modern global economy.}
{Đây là một tính từ dùng để chỉ các thực thể đa quốc gia, thường là các tập đoàn lớn.}

\vocabentry{Interdependence}
{/ˌɪntərdɪˈpendəns/ (n)}
{The dependence of two or more people or things on each other.}
{Global trade has created a state of interdependence where countries rely on each other for vital goods.}
{Đây là một danh từ dùng để chỉ sự phụ thuộc lẫn nhau giữa các quốc gia về mặt tài nguyên và kinh tế.}

\vocabentry{Liberalization}
{/ˌlɪbrələˈzeɪʃn/ (n)}
{The removal or reduction of restrictions or barriers on the free exchange of goods between nations.}
{Trade liberalization has encouraged many developing nations to open their markets to foreign investment.}
{Đây là một danh từ dùng để chỉ sự tự do hóa, đặc biệt là tự do hóa thương mại.}

\vocabentry{Outsourcing}
{/ˈaʊtsɔːrsɪŋ/ (n)}
{Obtain (goods or a service) from an outside or foreign supplier, especially in place of an internal source.}
{Many tech companies use outsourcing to call centers in India to save on operational costs.}
{Đây là một danh từ dùng để chỉ việc thuê ngoài nhân lực hoặc dịch vụ từ nước ngoài để giảm chi phí.}

\vocabentry{Homogenization}
{/həˌmɑːdʒənəˈzeɪʃn/ (n)}
{The process of making things uniform or similar.}
{Critics of globalization fear the cultural homogenization of the world, where local traditions are lost.}
{Đây là một danh từ dùng để chỉ sự đồng nhất hóa, thường nói về việc văn hóa thế giới trở nên giống nhau do ảnh hưởng của phương Tây.}

\vocabentry{Transnational}
{/trænzˈnæʃnəl/ (adj)}
{Extending or operating across national boundaries.}
{Transnational crime requires international cooperation between police forces to solve.}
{Đây là một tính từ dùng để chỉ các hoạt động xuyên quốc gia.}

\vocabentry{Tariff}
{/ˈtærɪf/ (n)}
{A tax or duty to be paid on a particular class of imports or exports.}
{The government imposed high tariffs on imported steel to protect local manufacturers.}
{Đây là một danh từ dùng để chỉ thuế quan, một công cụ bảo hộ mậu dịch.}

\vocabentry{Deregulation}
{/ˌdiːˌreɡjuˈleɪʃn/ (n)}
{The removal of regulations or restrictions, especially in a particular industry.}
{Deregulation of the banking sector led to an increase in cross-border financial transactions.}
{Đây là một danh từ dùng để chỉ sự bãi bỏ các quy định hạn chế trong kinh doanh và tài chính.}

\vocabentry{Cultural exchange}
{/ˈkʌltʃərəl ɪksˈtʃeɪndʒ/ (n)}
{An exchange of students, artists, etc., between two countries to promote mutual understanding.}
{International festivals are a great platform for cultural exchange and global friendship.}
{Đây là một cụm danh từ dùng để chỉ sự trao đổi văn hóa giữa các quốc gia.}

\vocabentry{Foreign Direct Investment}
{/ˈfɔːrən dɪˈrekt ɪnˈvestmənt/ (n)}
{An investment in the form of a controlling ownership in a business in one country by an entity based in another country. (Đây là một cụm danh từ dùng để chỉ vốn đầu tư trực tiếp nước ngoài (FDI).).}
{Foreign direct investment is a crucial driver of economic growth in many emerging markets.}
{}

\vocabentry{Migration}
{/maɪˈɡreɪʃn/ (n)}
{Movement of people to a new area or country in order to find work or better living conditions.}
{Global migration patterns are often driven by economic opportunities in developed nations.}
{Đây là một danh từ dùng để chỉ sự di cư, một hệ quả của toàn cầu hóa.}

\vocabentry{Commodity}
{/kəˈmɑːdəti/ (n)}
{A raw material or primary agricultural product that can be bought and sold.}
{Oil is perhaps the most important commodity in global trade.}
{Đây là một danh từ dùng để chỉ hàng hóa hoặc nguyên liệu thô giao dịch trên thị trường quốc tế.}

\vocabentry{Exploitation}
{/ˌeksplɔɪˈteɪʃn/ (n)}
{The action or fact of treating someone unfairly in order to benefit from their work.}
{There is a growing concern about the exploitation of workers in sweatshops in developing countries.}
{Đây là một danh từ dùng để chỉ sự khai thác hoặc bóc lột, thường nói về lao động giá rẻ ở nước nghèo.}

\vocabentry{Sustainability}
{/səˌsteɪnəˈbɪləti/ (n)}
{The ability to be maintained at a certain rate or level.}
{Globalization must be managed with an emphasis on environmental sustainability.}
{Đây là một danh từ dùng để chỉ tính bền vững, mục tiêu quan trọng trong phát triển toàn cầu.}

\vocabentry{Standardization}
{/ˌstændərdəˈzeɪʃn/ (n)}
{The process of making something conform to a standard.}
{The standardization of shipping containers has revolutionized international logistics.}
{Đây là một danh từ dùng để chỉ sự tiêu chuẩn hóa các sản phẩm và dịch vụ trên toàn thế giới.}

\vocabentry{Borderless}
{/ˈbɔːrdərləs/ (adj)}
{Not having borders; permitting free movement of people or goods.}
{The internet has created a borderless world for the exchange of ideas and information.}
{Đây là một tính từ dùng để chỉ một thế giới "không biên giới" nhờ vào internet và thương mại tự do.}

\vocabentry{Imperialism}
{/ɪmˈpɪriəlɪzəm/ (n)}
{A policy of extending a country's power and influence through diplomacy or military force.}
{Some view the dominance of American movies as a form of cultural imperialism.}
{Đây là một danh từ dùng để chỉ chủ nghĩa đế quốc, đôi khi dùng để chỉ "cultural imperialism".}

\vocabentry{Emerging market}
{/ɪˈmɜːrdʒɪŋ ˈmɑːrkɪt/ (n)}
{A nation's economy that is progressing toward becoming advanced.}
{Investors are looking for high returns in emerging markets like Vietnam and Brazil.}
{Đây là một cụm danh từ dùng để chỉ thị trường mới nổi, các nước đang phát triển nhanh.}

\vocabentry{Consolidated}
{/kənˈsɑːlɪdeɪtɪd/ (adj)}
{Combined into a single more effective or coherent whole.}
{The media industry is being consolidated into a few giant transnational corporations.}
{Đây là một tính từ chỉ sự hợp nhất các nguồn lực hoặc công ty toàn cầu.}

\vocabentry{Connectivity}
{/ˌkɑːnekˈtɪvəti/ (n)}
{The state or quality of being connected or connective.}
{High-speed internet has greatly improved connectivity between remote parts of the world.}
{Đây là một danh từ dùng để chỉ khả năng kết nối toàn cầu thông qua công nghệ.}

\vocabentry{Supply chain}
{/ˈsʌplaɪ tʃeɪn/ (n)}
{The sequence of processes involved in the production and distribution of a commodity.}
{A disruption in the global supply chain can cause shortages of goods everywhere.}
{Đây là một cụm danh từ dùng để chỉ chuỗi cung ứng toàn cầu.}

\vocabentry{Offshoring}
{/ˌɔːfˈʃɔːrɪŋ/ (n)}
{The practice of basing some of a company's processes or services overseas.}
{Offshoring production can increase profits but may lead to job losses in the home country.}
{Đây là một danh từ dùng để chỉ việc dời các hoạt động kinh doanh ra nước ngoài để tận dụng chi phí thấp.}

\vocabentry{Westernization}
{/ˌwestərnəˈzeɪʃn/ (n)}
{The adoption of the practices and culture of western Europe and North America.}
{Many traditional societies are experiencing rapid westernization due to global media influence.}
{Đây là một danh từ dùng để chỉ sự phương Tây hóa.}

\vocabentry{Free trade}
{/ˌfriː ˈtreɪd/ (n)}
{International trade left to its natural course without tariffs, quotas, or other restrictions.}
{Free trade agreements aim to lower costs for consumers and increase competition.}
{Đây là một cụm danh từ dùng để chỉ thương mại tự do.}

\vocabentry{Privatization}
{/ˌpraɪvətəˈzeɪʃn/ (n)}
{The transfer of a business, industry, or service from public to private ownership and control.}
{Privatization of state-owned enterprises is often a requirement for global financial aid.}
{Đây là một danh từ dùng để chỉ sự tư nhân hóa các dịch vụ công.}

\vocabentry{Infrastructure}
{/ˈɪnfrəstrʌktʃər/ (n)}
{The basic physical and organizational structures and facilities needed for the operation of a society.}
{Developing world-class infrastructure is essential for attracting multinational investment.}
{Đây là một danh từ dùng để chỉ cơ sở hạ tầng phục vụ kết nối toàn cầu.}

\vocabentry{Monopoly}
{/əˈmɑːpəli/ (n)}
{The exclusive possession or control of the supply of or trade in a commodity or service.}
{Global anti-trust laws aim to prevent tech giants from holding a monopoly on information.}
{Đây là một danh từ dùng để chỉ sự độc quyền của các tập đoàn lớn.}

\vocabentry{Remittance}
{/rɪˈmɪtns/ (n)}
{A sum of money sent, especially by mail, in payment for goods or services or as a gift.}
{Remittances from citizens working abroad are a vital source of income for many families.}
{Đây là một danh từ dùng để chỉ kiều hối, tiền người lao động di cư gửi về nước.}

\vocabentry{Affluence}
{/ˈæfluəns/ (n)}
{The state of having a great deal of money; wealth.}
{Increased global trade has brought newfound affluence to many middle-class families in Asia.}
{Đây là một danh từ dùng để chỉ sự giàu có, sung túc mang lại nhờ thương mại toàn cầu.}

\vocabentry{Human rights}
{/ˌhjuːmən ˈraɪts/ (n)}
{Rights which are believed to belong justifiably to every person.}
{International organizations work to ensure that human rights are respected in all global trade deals.}
{Đây là một cụm danh từ dùng để chỉ quyền con người, một tiêu chuẩn toàn cầu.}

\vocabentry{Labor market}
{/ˈleɪbər ˈmɑːrkɪt/ (n)}
{The supply of people in a particular country or area who are able and willing to work.}
{Globalization has forced workers to compete in a global labor market.}
{Đây là một cụm danh từ dùng để chỉ thị trường lao động toàn cầu.}

\vocabentry{Digital divide}
{/ˌdɪdʒɪtl dɪˈvaɪd/ (n)}
{The gulf between those who have ready access to computers and the internet, and those who do not.}
{The digital divide remains a significant barrier to equal participation in the global economy.}
{Đây là một cụm danh từ dùng để chỉ khoảng cách công nghệ giữa các quốc gia.}

\vocabentry{Prosperity}
{/prɑːˈsperəti/ (n)}
{The state of being prosperous; wealth and success.}
{The ultimate goal of economic globalization is to bring prosperity to all corners of the globe.}
{Đây là một danh từ dùng để chỉ sự thịnh vượng.}

\vocabentry{Intercultural}
{/ˌɪntərˈkʌltʃərəl/ (adj)}
{Relating to or involving different cultures.}
{Education should promote intercultural competence in a globalized world.}
{Đây là một tính từ dùng để chỉ các mối quan hệ hoặc sự hiểu biết giữa các nền văn hóa.}

\vocabentry{Global village}
{/ˌɡloʊbl ˈvɪlɪdʒ/ (n)}
{The world considered as a single community linked by telecommunications.}
{Telecommunications have effectively turned the world into a global village.}
{Đây là một cụm danh từ dùng để chỉ "ngôi làng toàn cầu".}

\vocabentry{Expatriate}
{/ˌeksˈpeɪtriət/ (n)}
{A person who lives outside their native country.}
{There is a large expatriate community of software engineers in Singapore.}
{Đây là một danh từ dùng để chỉ người sống và làm việc ở nước ngoài.}

\vocabentry{Decentralization}
{/ˌdiːˌsentrələˈzeɪʃn/ (n)}
{The transfer of control of an activity or organization to several local offices or authorities.}
{Global companies often use decentralization to adapt to local market needs.}
{Đây là một danh từ dùng để chỉ sự phân quyền, tản quyền.}

\vocabentry{Brain drain}
{/ˈbreɪn dreɪn/ (n)}
{The emigration of highly trained or intelligent people from a particular country.}
{Developing nations often suffer from brain drain as their best doctors move to the West.}
{Đây là một cụm danh từ dùng để chỉ hiện tượng chảy máu chất xám.}

\vocabentry{Brain gain}
{/ˈbreɪn ɡeɪn/ (n)}
{An increase in the number of highly trained or intelligent people in a particular country.}
{The government introduced new visas to promote brain gain in the tech sector.}
{Đây là một cụm danh từ dùng để chỉ hiện tượng thu hút chất xám.}

\vocabentry{Standard of living}
{/ˌstændərd əv ˈlɪvɪŋ/ (n)}
{The degree of wealth and material comfort available to a person or community.}
{Globalization has helped raise the standard of living for millions of people in poverty.}
{Đây là một cụm danh từ dùng để chỉ mức sống.}

\vocabentry{Inequality}
{/ˌɪnɪˈkwɑːləti/ (n)}
{Difference in size, degree, circumstances, etc.; lack of equality.}
{Economic globalization is often criticized for increasing income inequality.}
{Đây là một danh từ dùng để chỉ sự bất bình đẳng giàu nghèo giữa các quốc gia.}

\vocabentry{Cooperation}
{/koʊˌɑːpəˈreɪʃn/ (n)}
{The process of working together to the same end.}
{Global cooperation is necessary to combat climate change effectively.}
{Đây là một danh từ dùng để chỉ sự hợp tác quốc tế.}

\vocabentry{Regulation}
{/ˌreɡjuˈleɪʃn/ (n)}
{A rule or directive made and maintained by an authority.}
{New environmental regulations are being applied to international shipping companies.}
{Đây là một danh từ dùng để chỉ các quy định pháp lý quản lý thương mại toàn cầu.}

\vocabentry{Multiculturalism}
{/ˌmʌltiˈkʌltʃərəlɪzəm/ (n)}
{The presence of, or support for the presence of, several distinct cultural or ethnic groups within a society.}
{Big cities like New York are famous for their vibrant multiculturalism.}
{Đây là một danh từ dùng để chỉ chủ nghĩa đa văn hóa.}

\vocabentry{Linguistic}
{/lɪŋˈɡwɪstɪk/ (adj)}
{Relating to language or linguistics.}
{The linguistic dominance of English is a major feature of globalization.}
{Đây là một tính từ dùng để chỉ các yếu tố ngôn ngữ trong toàn cầu hóa.}

\vocabentry{Cosmopolitan}
{/ˌkɑːzməˈpɑːlɪtən/ (adj)}
{Familiar with and at ease in many different countries and cultures.}
{Studying abroad can help students develop a more cosmopolitan outlook on life.}
{Đây là một tính từ dùng để chỉ một người hoặc một nơi mang tầm vóc quốc tế, am hiểu nhiều văn hóa.}

\vocabentry{Diplomacy}
{/dɪˈploʊməsi/ (n)}
{The profession, activity, or skill of managing international relations.}
{Digital diplomacy is becoming increasingly important in international relations.}
{Đây là một danh từ dùng để chỉ ngành ngoại giao.}

\vocabentry{Sanction}
{/ˈæŋkʃn/ (n)}
{A threatened penalty for disobeying a law or rule.}
{The international community imposed economic sanctions to pressure the government.}
{Đây là một danh từ dùng để chỉ sự trừng phạt hoặc cấm vận kinh tế quốc tế.}

\vocabentry{Universal}
{/ˌjuːnɪˈvɜːrsl/ (adj)}
{Of, affecting, or done by all people or things in the world or in a particular group.}
{There are universal human rights that should be protected regardless of nationality.}
{Đây là một tính từ dùng để chỉ tính phổ quát, toàn cầu.}

\vocabentry{Logistics}
{/ləˈdʒɪstɪks/ (n)}
{The detailed coordination of a complex operation involving many people, facilities, or supplies.}
{Efficient logistics are the backbone of the global e-commerce industry.}
{Đây là một danh từ dùng để chỉ ngành hậu cần, vận tải hàng hóa quốc tế.}

\vocabentry{E-commerce}
{/ˈiː kɑːmɜːrs/ (n)}
{Commercial transactions conducted electronically on the internet.}
{E-commerce allows small businesses to reach customers all over the world.}
{Đây là một danh từ dùng để chỉ thương mại điện tử.}

\vocabentry{Developing nation}
{/dɪˈveləpɪŋ ˈneɪʃn/ (n)}
{A country with a less developed industrial base and a low Human Development Index.}
{Globalization provides both opportunities and risks for a developing nation.}
{Đây là một cụm danh từ dùng để chỉ các quốc gia đang phát triển.}

\vocabentry{Developed nation}
{/dɪˈveləpt ˈneɪʃn/ (n)}
{A sovereign state that has a high quality of life, developed economy, and advanced technological infrastructure.}
{Developed nations are often the primary beneficiaries of global trade.}
{Đây là một cụm danh từ dùng để chỉ các quốc gia phát triển.}

\vocabentry{Exploitative}
{/ɪkˈsplɔɪtətɪv/ (adj)}
{Making use of a situation or person in an unfair or selfish way.}
{Consumer groups are boycotting companies with exploitative labor practices.}
{Đây là một tính từ dùng để chỉ tính chất bóc lột.}

\vocabentry{Harmonization}
{/ˌhɑːrmənəˈzeɪʃn/ (n)}
{The process of making different systems or rules similar or compatible.}
{The harmonization of tax laws would help reduce global tax evasion.}
{Đây là một danh từ dùng để chỉ sự hài hòa hóa các luật lệ thương mại giữa các nước.}

\vocabentry{Fragile}
{/ˈfrædʒl/ (adj)}
{(Of an object) easily broken or damaged.}
{Small island nations have fragile economies that are highly sensitive to global market changes.}
{Trong toàn cầu hóa, dùng để chỉ sự mong manh của các nền kinh tế nhỏ.}

\vocabentry{Biodiversity}
{/ˌbaɪoʊdaɪˈvɜːrsəti/ (n)}
{The variety of plant and animal life in the world or in a particular habitat.}
{Globalization can threaten biodiversity through the introduction of invasive species.}
{Đây là một danh từ dùng để chỉ đa dạng sinh học, một vấn đề môi trường toàn cầu.}

\vocabentry{Climate change}
{/ˈklaɪmət tʃeɪndʒ/ (n)}
{A change in global or regional climate patterns.}
{Climate change is a global problem that requires a global solution.}
{Đây là một cụm danh từ dùng để chỉ biến đổi khí hậu.}

\vocabentry{Outflow}
{/ˈaʊtfloʊ/ (n)}
{A large amount of something that moves or is transferred out of a place.}
{The sudden outflow of capital caused a financial crisis in the country.}
{Đây là một danh từ dùng để chỉ sự chảy ra ngoài của vốn hoặc nhân lực.}

\vocabentry{Inflow}
{/ˈɪnfloʊ/ (n)}
{A large amount of something that moves or is transferred into a place.}
{The country has seen a massive inflow of foreign investment this year.}
{Đây là một danh từ dùng để chỉ sự chảy vào của vốn đầu tư.}

\vocabentry{Regionalism}
{/ˈriːdʒənəlɪzəm/ (n)}
{The theory or practice of regional rather than central systems of administration or economic, cultural, or political affiliation.}
{Regionalism can sometimes act as a stepping stone to full globalization.}
{Đây là một danh từ dùng để chỉ chủ nghĩa khu vực.}

\vocabentry{Protectionism}
{/prəˈtekʃənɪzəm/ (n)}
{The theory or practice of shielding a country's domestic industries from foreign competition.}
{Protectionism can save domestic jobs in the short term but may hurt consumers in the long run.}
{Đây là một danh từ dùng để chỉ chủ nghĩa bảo hộ mậu dịch.}

\vocabentry{Quota}
{/ˈkwoʊtə/ (n)}
{A fixed share of something that a person or group is entitled to receive or is bound to contribute.}
{The government has placed a quota on the amount of foreign textiles that can be imported.}
{Trong thương mại, đây là hạn ngạch nhập khẩu.}

\vocabentry{Subsidy}
{/ˈsʌbsədi/ (n)}
{A sum of money granted by the government or a public body to assist an industry or business.}
{Farm subsidies in wealthy nations make it difficult for farmers in poor countries to compete.}
{Đây là một danh từ dùng để chỉ sự trợ cấp của chính phủ.}

\vocabentry{World Trade Organization}
{/wɜːrld treɪd ˌɔːrɡənəˈzeɪʃn/ (n)}
{An international organization that regulates and facilitates international trade. (Đây là một danh từ riêng dùng để chỉ Tổ chức Thương mại Thế giới (WTO).).}
{The World Trade Organization helps resolve trade disputes between member countries.}
{}

\vocabentry{International Monetary Fund}
{/ˌɪntərˈnæʃnəl ˈmɑːnɪteri fʌnd/ (n)}
{An international organization that works to foster global monetary cooperation. (Đây là một danh từ riêng dùng để chỉ Quỹ Tiền tệ Quốc tế (IMF).).}
{The International Monetary Fund provided a loan to the country during the economic crisis.}
{}

\vocabentry{World Bank}
{/wɜːrld bæŋk/ (n)}
{An international financial institution that provides loans and grants to the governments of low- and middle-income countries.}
{The World Bank funds many infrastructure projects in developing nations.}
{Đây là một danh từ riêng dùng để chỉ Ngân hàng Thế giới.}

\vocabentry{Free movement}
{/ˌfriː ˈmuːvmənt/ (n)}
{The right of citizens of a state to travel, reside, and work in any part of the state or union.}
{Free movement of labor is a core principle of the European Union.}
{Đây là một cụm danh từ dùng để chỉ sự tự do di chuyển.}

\vocabentry{Indigenous}
{/ɪnˈdɪdʒənəs/ (adj)}
{Originating or occurring naturally in a particular place; native.}
{Global companies must respect the land rights of indigenous peoples.}
{Đây là một tính từ dùng để chỉ các dân tộc hoặc văn hóa bản địa.}

\vocabentry{Authenticity}
{/ˌɔːθenˈtɪsəti/ (n)}
{The quality of being authentic.}
{Many tourists travel to remote areas in search of cultural authenticity.}
{Trong văn hóa, dùng để chỉ tính nguyên bản của văn hóa địa phương.}

\vocabentry{Acculturation}
{/əˌkʌltʃəˈreɪʃn/ (n)}
{The process of social, psychological, and cultural change that stems from the balancing of two cultures.}
{Acculturation can lead to the development of new, hybrid cultural forms.}
{Đây là một danh từ dùng để chỉ sự tiếp nhận văn hóa.}

\vocabentry{Hybridization}
{/ˌhaɪbrɪdəˈzeɪʃn/ (n)}
{The process of mixing different cultures or styles.}
{Global food chains often practice cultural hybridization by offering local menu items.}
{Đây là một danh từ dùng để chỉ sự lai tạp hoặc kết hợp văn hóa.}

\vocabentry{Digital economy}
{/ˌdɪdʒɪtl ɪˈkɑːnəmi/ (n)}
{An economy that is based on digital computing technologies.}
{The digital economy has allowed people to work for global companies from their own homes.}
{Đây là một cụm danh từ dùng để chỉ nền kinh tế số.}

\vocabentry{Labor standards}
{/ˈleɪbər ˈstændərdz/ (n)}
{The rules that govern how people work.}
{Improving global labor standards is necessary to prevent a "race to the bottom" in wages.}
{Đây là một cụm danh từ dùng để chỉ các tiêu chuẩn lao động quốc tế.}

\vocabentry{Race to the bottom}
{/reɪs tuː ðə ˈbɑːtəm/ (idiom)}
{A situation in which companies or countries compete to provide the lowest cost or lowest regulation.}
{Unregulated globalization can lead to a race to the bottom in environmental protection.}
{Đây là một thành ngữ dùng để chỉ sự cạnh tranh xuống đáy về giá cả hoặc tiêu chuẩn.}

\vocabentry{Intellectual property}
{/ˌɪntəˈlektʃuəl ˈprɑːpərti/ (n)}
{A work or invention that is the result of creativity.}
{Protection of intellectual property is a major issue in international trade negotiations.}
{Đây là một cụm danh từ dùng để chỉ sở hữu trí tuệ.}

\vocabentry{Patent}
{/ˈpætnt/ (n)}
{A government authority or license conferring a right or title for a set period.}
{Global pharmaceutical companies hold patents that can make medicines very expensive.}
{Đây là một danh từ dùng để chỉ bằng sáng chế.}

\vocabentry{Royalty}
{/ˈrɔɪəlti/ (n)}
{A sum of money paid to a patentee for the use of a patent.}
{The artist receives royalties whenever her music is played on international radio.}
{Đây là một danh từ dùng để chỉ tiền bản quyền.}

\vocabentry{Ethical}
{/ˈeθɪkl/ (adj)}
{Relating to moral principles.}
{Many consumers now prefer to buy products from companies with ethical supply chains.}
{Đây là một tính từ dùng để chỉ tính đạo đức trong kinh doanh toàn cầu.}

\vocabentry{Fair trade}
{/ˌfer ˈtreɪd/ (n)}
{Trade between companies in developed countries and producers in developing countries in which fair prices are paid to the producers.}
{Buying fair trade coffee ensures that farmers receive a living wage.}
{Đây là một cụm danh từ dùng để chỉ thương mại công bằng.}

\vocabentry{Social justice}
{/ˌsoʊʃl ˈdʒʌstɪs/ (n)}
{Justice in terms of the distribution of wealth, opportunities, and privileges within a society.}
{Global movements for social justice aim to reduce the gap between the rich and the poor.}
{Đây là một cụm danh từ dùng để chỉ công bằng xã hội.}

\vocabentry{Exploit}
{/ɪkˈsplɔɪt/ (v)}
{Make full use of and derive benefit from (a resource).}
{Companies seek to exploit new market opportunities in Southeast Asia.}
{Đây là một động từ mang nghĩa tích cực là tận dụng, hoặc tiêu cực là bóc lột.}

\vocabentry{Vulnerable}
{/ˈvʌlnərəbl/ (adj)}
{Susceptible to physical or emotional attack or harm.}
{Small farmers are often vulnerable to sudden drops in global commodity prices.}
{Tính từ dùng để chỉ sự dễ bị tổn thương của các nước nghèo.}

\vocabentry{Resilience}
{/rɪˈzɪliəns/ (n)}
{The capacity to recover quickly from difficulties.}
{Diversifying the economy can help build resilience against global financial shocks.}
{Đây là một danh từ dùng để chỉ khả năng phục hồi của nền kinh tế.}

\vocabentry{Diversification}
{/daɪˌvɜːrsɪfɪˈkeɪʃn/ (n)}
{The process of a business enlarging or varying its range of products or field of operation.}
{Market diversification is essential for reducing a country's dependence on a single export.}
{Đây là một danh từ dùng để chỉ sự đa dạng hóa sản phẩm hoặc thị trường.}

\vocabentry{Competitive advantage}
{/kəmˈpetətɪv ədˈvæntɪdʒ/ (n)}
{A condition or circumstance that puts a company in a favorable or superior business position.}
{A highly skilled workforce is a significant competitive advantage in the global economy.}
{Đây là một cụm danh từ dùng để chỉ lợi thế cạnh tranh.}

\vocabentry{Comparative advantage}
{/kəmˈpærətɪv ədˈvæntɪdʒ/ (n)}
{The ability of an individual or group to carry out a particular economic activity more efficiently than another activity.}
{According to the theory of comparative advantage, countries should specialize in what they do best.}
{Đây là một cụm danh từ dùng để chỉ lợi thế so sánh.}

\vocabentry{Specialize}
{/ˈspeʃəlaɪz/ (v)}
{Concentrate on and become expert in a particular subject or skill.}
{Some countries specialize in high-tech manufacturing, while others focus on tourism.}
{Đây là một động từ dùng để chỉ sự chuyên môn hóa sản xuất.}

\vocabentry{Economy of scale}
{/ɪˈkɑːnəmi əv ˈskeɪl/ (n)}
{A proportionate saving in costs gained by an increased level of production.}
{Multinational companies benefit from economies of scale that small businesses cannot match.}
{Đây là một cụm danh từ dùng để chỉ lợi thế kinh tế nhờ quy mô.}

\vocabentry{Market share}
{/ˈmɑːrkɪt ʃer/ (n)}
{The portion of a market controlled by a particular company or product.}
{The tech company is aggressively expanding its market share in Europe.}
{Đây là một cụm danh từ dùng để chỉ thị phần.}

\vocabentry{Penetrate}
{/ˈpenətreɪt/ (v)}
{Go into or through (something); in business, to enter a market.}
{It can be difficult for foreign brands to penetrate the Japanese market.}
{Đây là một động từ dùng để chỉ việc thâm nhập thị trường.}

\vocabentry{Saturation}
{/ˌsætʃəˈreɪʃn/ (n)}
{The state or process that happens when no more of something can be accepted or added.}
{The smartphone market in developed nations has reached a point of saturation.}
{Đây là một danh từ dùng để chỉ sự bão hòa thị trường.}

\vocabentry{Consumption}
{/kənˈsʌmpʃn/ (n)}
{The using up of a resource.}
{Global patterns of consumption are changing as people become more environmentally conscious.}
{Đây là một danh từ dùng để chỉ sự tiêu dùng.}

\vocabentry{Middle class}
{/ˌmɪdl ˈklæs/ (n)}
{The social group between the upper and working classes.}
{The growth of the global middle class is creating huge new markets for luxury goods.}
{Đây là một cụm danh từ dùng để chỉ tầng lớp trung lưu, động lực tiêu dùng toàn cầu.}

\vocabentry{Urban}
{/ˈɜːrbən/ (adj)}
{In, relating to, or characteristic of a town or city.}
{Globalization has accelerated urban development across the world.}
{Đây là một tính từ dùng để chỉ các khu vực đô thị.}

\vocabentry{Rural}
{/ˈrʊrəl/ (adj)}
{In, relating to, or characteristic of the countryside rather than the town.}
{Many people in rural areas feel left behind by the benefits of globalization.}
{Đây là một tính từ dùng để chỉ các khu vực nông thôn.}

\vocabentry{Megacity}
{/ˈmeɡəsɪti/ (n)}
{A very large city, typically one with a population of over ten million people.}
{Tokyo and New York are examples of megacities that act as global hubs.}
{Đây là một danh từ dùng để chỉ các siêu đô thị toàn cầu.}

\vocabentry{Hub}
{/hʌb/ (n)}
{The effective center of an activity, region, or network.}
{Singapore is a major global hub for finance and shipping.}
{Đây là một danh từ dùng để chỉ trung tâm hoặc đầu mối toàn cầu.}

\vocabentry{Interconnection}
{/ˌɪntərkəˈnekʃn/ (n)}
{A mutual connection between two or more things.}
{The interconnection of global financial markets means that a crisis in one country can spread quickly.}
{Đây là một danh từ dùng để chỉ sự liên kết lẫn nhau.}

\vocabentry{Flow}
{/floʊ/ (n/v)}
{The action or fact of moving along in a steady, continuous stream.}
{The free flow of information is essential for a functioning global democracy.}
{Trong toàn cầu hóa, dùng để chỉ luồng vốn, hàng hóa hoặc thông tin.}

\vocabentry{Capital}
{/ˈkæpɪtl/ (n)}
{Wealth in the form of money or other assets owned by a person or organization.}
{Capital can move across borders at the click of a button in today's world.}
{Đây là một danh từ dùng để chỉ vốn đầu tư.}

\vocabentry{Labor}
{/ˈleɪbər/ (n)}
{Work, especially hard physical work.}
{Companies often move production to countries where labor is cheaper.}
{Đây là một danh từ dùng để chỉ nhân công hoặc lực lượng lao động.}

\vocabentry{Skilled labor}
{/ˈskɪld ˈleɪbər/ (n)}
{Labor that requires specialized training or skills.}
{There is a global shortage of skilled labor in the renewable energy sector.}
{Đây là một cụm danh từ dùng để chỉ lao động có tay nghề cao.}

\vocabentry{Unskilled labor}
{/ˌʌnˈskɪld ˈleɪbər/ (n)}
{Labor that does not require specialized training.}
{Globalization has made life difficult for unskilled labor in high-wage countries.}
{Đây là một cụm danh từ dùng để chỉ lao động phổ thông.}

\vocabentry{Workforce}
{/ˈwɜːrkfɔːrs/ (n)}
{The people engaged in or available for work, either in a country or area or in a particular company.}
{The company has a diverse workforce representing over twenty different nationalities.}
{Đây là một danh từ dùng để chỉ lực lượng lao động.}

\vocabentry{Minimum wage}
{/ˌmɪnɪməm ˈweɪdʒ/ (n)}
{The lowest wage permitted by law or by a special agreement.}
{Some activists argue for a global minimum wage to protect vulnerable workers.}
{Đây là một cụm danh từ dùng để chỉ mức lương tối thiểu.}

\vocabentry{Sweatshop}
{/ˈswetʃɑːp/ (n)}
{A factory or workshop, especially in the clothing industry, where manual workers are employed at very low wages for long hours and under poor conditions.}
{Many global brands have been accused of using sweatshops in their supply chains.}
{Đây là một danh từ dùng để chỉ xưởng bóc lột sức lao động.}

\vocabentry{Trade union}
{/ˈtreɪd ˌjuːniən/ (n)}
{An organized association of workers formed to protect and further their rights and interests.}
{Trade unions are struggling to adapt to the reality of the global labor market.}
{Đây là một cụm danh từ dùng để chỉ công đoàn lao động.}

\vocabentry{Boycott}
{/ˈbɔɪkɑːt/ (v/n)}
{Withdraw from commercial or social relations with (a country, organization, or person) as a punishment or protest.}
{Consumers organized a boycott of the company to protest its environmental record.}
{Đây là một động từ/danh từ dùng để chỉ sự tẩy chay.}

\vocabentry{Activism}
{/ˈæktɪvɪzəm/ (n)}
{The policy or action of using vigorous campaigning to bring about political or social change.}
{Digital activism has allowed global movements to organize more effectively.}
{Đây là một danh từ dùng để chỉ chủ nghĩa hoạt động, phong trào xã hội.}

\vocabentry{NGO}
{/ˌen dʒiː ˈoʊ/ (n)}
{Non-governmental organization; a non-profit organization that operates independently of any government.}
{Many NGOs work on global issues like poverty, education, and health.}
{Đây là một danh từ viết tắt dùng để chỉ các tổ chức phi chính phủ.}

\vocabentry{Civil society}
{/ˌsɪvl səˈsaɪəti/ (n)}
{Society considered as a community of citizens linked by common interests and collective activity.}
{A strong civil society is essential for holding global corporations accountable.}
{Đây là một cụm danh từ dùng để chỉ xã hội dân sự.}

\vocabentry{Advocacy}
{/ˈædvəkəsi/ (n)}
{Public support for or recommendation of a particular cause or policy.}
{The group's advocacy for fair trade has led to significant policy changes.}
{Đây là một danh từ dùng để chỉ sự vận động, ủng hộ một chính sách nào đó.}

\vocabentry{Accountability}
{/əˌkaʊntəˈbɪləti/ (n)}
{The fact or condition of being accountable; responsibility.}
{There are calls for greater accountability for transnational corporations.}
{Đây là một danh từ dùng để chỉ trách nhiệm giải trình.}

\vocabentry{Transparency}
{/trænsˈpærənsi/ (n)}
{The condition of being transparent; in business, being open about operations and finances.}
{Transparency in the supply chain is vital for ethical consumerism.}
{Đây là một danh từ dùng để chỉ sự minh bạch.}

\vocabentry{Corruption}
{/kəˈrʌpʃn/ (n)}
{Dishonest or fraudulent conduct by those in power.}
{Corruption can be a major barrier to economic development in a globalized world.}
{Đây là một danh từ dùng để chỉ sự tham nhũng.}

\vocabentry{Bribe}
{/braɪb/ (n/v)}
{Dishonestly persuade (someone) to act in one's favor by a gift of money or other inducement.}
{International laws aim to prevent companies from using bribes to win foreign contracts.}
{Đây là một danh từ/động từ dùng để chỉ sự hối lộ.}

\vocabentry{Lobby}
{/ˈlɑːbi/ (v/n)}
{Seek to influence (a politician or public official) on an issue.}
{Tech giants spend millions each year to lobby governments on data privacy laws.}
{Đây là một động từ/danh từ dùng để chỉ hoạt động vận động hành lang.}

\vocabentry{Neoliberalism}
{/ˌniːoʊˈlɪbrəlɪzəm/ (n)}
{A political approach that favors free-market capitalism, deregulation, and reduction in government spending.}
{Neoliberalism has been the dominant economic ideology behind globalization since the 1980s.}
{Đây là một danh từ dùng để chỉ chủ nghĩa tân tự do.}

\vocabentry{Capitalism}
{/ˈæpɪtəlɪzəm/ (n)}
{An economic and political system in which a country's trade and industry are controlled by private owners for profit.}
{Global capitalism has driven innovation but also created significant wealth gaps.}
{Đây là một danh từ dùng để chỉ chủ nghĩa tư bản.}

\vocabentry{Socialism}
{/ˈsoʊʃəlɪzəm/ (n)}
{A political and economic theory of social organization which advocates that the means of production should be owned or regulated by the community.}
{Some countries follow a model of "democratic socialism" with strong public services.}
{Đây là một danh từ dùng để chỉ chủ nghĩa xã hội.}

\vocabentry{Democracy}
{/dɪˈmɑːkrəsi/ (n)}
{A system of government by the whole population or all the eligible members of a state.}
{The spread of democracy is often seen as a positive effect of globalization.}
{Đây là một danh từ dùng để chỉ chế độ dân chủ.}

\vocabentry{Governance}
{/ˈɡʌvərnəns/ (n)}
{The action or manner of governing.}
{Good global governance is needed to manage the risks of financial crises.}
{Đây là một danh từ dùng để chỉ sự quản trị, điều hành.}

\vocabentry{Institution}
{/ˌɪnstɪˈtuːʃn/ (n)}
{An organization founded for a religious, educational, professional, or social purpose.}
{International institutions like the UN help maintain global peace and security.}
{Đây là một danh từ dùng để chỉ các tổ chức hoặc định chế quốc tế.}

\vocabentry{Multilateral}
{/ˌmʌltiˈlætərəl/ (adj)}
{Agreed upon or participated in by three or more parties, especially the governments of different countries.}
{Multilateral trade agreements are more complex but offer wider benefits than bilateral ones.}
{Đây là một tính từ dùng để chỉ tính đa phương.}

\vocabentry{Bilateral}
{/ˌbaɪˈlætərəl/ (adj)}
{Involving two parties, especially two countries.}
{The two nations signed a bilateral agreement to increase tourism.}
{Đây là một tính từ dùng để chỉ tính song phương.}

\vocabentry{Treaty}
{/ˈtriːti/ (n)}
{A formally concluded and ratified agreement between countries.}
{The climate treaty aims to reduce global carbon emissions by 50\%.}
{Đây là một danh từ dùng để chỉ hiệp ước quốc tế.}

\vocabentry{Accord}
{/əˈkɔːrd/ (n)}
{An official agreement or treaty.}
{The Paris Accord is a landmark agreement on global climate action.}
{Đây là một danh từ đồng nghĩa với hiệp định hoặc thỏa thuận.}

\vocabentry{Protocol}
{/ˈproʊtəkɑːl/ (n)}
{The official procedure or system of rules governing affairs of state or diplomatic occasions.}
{Diplomatic protocol must be strictly followed during international summits.}
{Đây là một danh từ dùng để chỉ nghi thức hoặc nghị định thư.}

\vocabentry{Summit}
{/ˈsʌmɪt/ (n)}
{A meeting between heads of government.}
{The G7 summit focused on global economic recovery after the pandemic.}
{Đây là một danh từ dùng để chỉ hội nghị thượng đỉnh.}

\vocabentry{Forum}
{/ˈfɔːrəm/ (n)}
{A place, meeting, or medium where ideas and views on a particular issue can be exchanged.}
{The World Economic Forum is held annually in Davos, Switzerland.}
{Đây là một danh từ dùng để chỉ diễn đàn quốc tế.}

\vocabentry{Delegation}
{/ˌdelɪˈɡeɪʃn/ (n)}
{A body of delegates or representatives; a group of people representing a country.}
{A delegation of business leaders traveled to China to explore new trade links.}
{Đây là một danh từ dùng để chỉ phái đoàn đại biểu.}

\vocabentry{Diplomat}
{/ˈdɪpləmæt/ (n)}
{An official representing a country abroad.}
{Diplomats work behind the scenes to resolve international conflicts.}
{Đây là một danh từ dùng để chỉ nhà ngoại giao.}

\vocabentry{Embassy}
{/ˈembəsi/ (n)}
{The official residence or offices of an ambassador.}
{The embassy provides assistance to its citizens traveling abroad.}
{Đây là một danh từ dùng để chỉ đại sứ quán.}

\vocabentry{Visa}
{/ˈviːzə/ (n)}
{An endorsement on a passport indicating that the holder is allowed to enter, leave, or stay for a specified period of time in a country.}
{Many countries have relaxed visa requirements to encourage global tourism.}
{Đây là một danh từ dùng để chỉ thị thực, một rào cản di chuyển.}

\vocabentry{Passport}
{/ˈpæspɔːrt/ (n)}
{An official document issued by a government, certifying the holder's identity and citizenship and entitling them to travel under its protection.}
{A powerful passport allows citizens to travel to many countries without a visa.}
{Đây là một danh từ dùng để chỉ hộ chiếu.}

\vocabentry{Citizen}
{/ˈsɪtɪzn/ (n)}
{A legally recognized subject or national of a state or commonwealth.}
{Many people now consider themselves "global citizens" rather than just citizens of one country.}
{Đây là một danh từ dùng để chỉ công dân.}

\vocabentry{Refugee}
{/ˌrefjuˈdʒiː/ (n)}
{A person who has been forced to leave their country in order to escape war, persecution, or natural disaster.}
{The international community has a responsibility to provide aid to refugees.}
{Đây là một danh từ dùng để chỉ người tị nạn.}

\vocabentry{Innovation}
{/ˌɪnəˈveɪʃn/ (n)}
{The action or process of innovating; a new method, idea, product, etc.}
{Technological innovation is the primary driver of global competition.}
{Đây là một danh từ dùng để chỉ sự đổi mới, cách tân.}

\vocabentry{Intellectual}
{/ˌɪntəˈlektʃuəl/ (adj)}
{Relating to the intellect; involving thought and reason.}
{The global exchange of intellectual ideas has accelerated scientific progress.}
{Đây là một tính từ dùng để chỉ các yếu tố trí tuệ, ví dụ: intellectual property.}

\vocabentry{Copyright}
{/ˈkɑːpiraɪt/ (n)}
{The exclusive legal right, given to an originator or an assignee to print, publish, perform, film, or record material.}
{Protecting copyright is difficult in the age of global digital sharing.}
{Đây là một danh từ dùng để chỉ bản quyền.}

\vocabentry{Trademark}
{/ˈtreɪdmɑːrk/ (n)}
{A symbol, word, or words legally registered or established by use as representing a company or product.}
{Global brands spend millions protecting their trademarks from counterfeiters.}
{Đây là một danh từ dùng để chỉ nhãn hiệu hàng hóa.}

\vocabentry{Counterfeit}
{/ˈkaʊntərfɪt/ (adj/n)}
{Made in exact imitation of something valuable or important with the intention to deceive or defraud.}
{The trade in counterfeit luxury goods is a multi-billion dollar global problem.}
{Đây là một tính từ/danh từ dùng để chỉ hàng giả.}

\vocabentry{Patent}
{/ˈpætnt/ (v)}
{Obtain a patent for (an invention).}
{Companies rush to patent new technologies to gain a global advantage.}
{Đây là một động từ dùng để chỉ việc đăng ký bằng sáng chế.}

\vocabentry{Research and Development}
{/rɪˈsɜːrtʃ ənd dɪˈveləpmənt/ (n)}
{Work directed toward the innovation, introduction, and improvement of products and processes. (Đây là một cụm danh từ dùng để chỉ nghiên cứu và phát triển (R\&D).).}
{Global tech firms invest heavily in research and development to stay ahead.}
{}

\vocabentry{Entrepreneur}
{/ˌɑːntrəprəˈnɜːr/ (n)}
{A person who organizes and operates a business or businesses, taking on greater than normal financial risks.}
{The internet has empowered entrepreneurs to start global businesses with very little capital.}
{Đây là một danh từ dùng để chỉ doanh nhân.}

\vocabentry{Startup}
{/ˈstɑːrtʌp/ (n)}
{A newly established business.}
{Many global startups are focused on solving environmental problems through technology.}
{Đây là một danh từ dùng để chỉ các công ty khởi nghiệp.}

\vocabentry{Venture capital}
{/ˈventʃər ˈkæpɪtl/ (n)}
{Capital invested in a project in which there is a substantial element of risk, typically a new or expanding business.}
{Venture capital is essential for the growth of global technology hubs.}
{Đây là một cụm danh từ dùng để chỉ vốn đầu tư mạo hiểm.}

\vocabentry{Monetary}
{/ˈmɑːnɪteri/ (adj)}
{Relating to money or currency.}
{Global monetary policy is often influenced by the decisions of the US Federal Reserve.}
{Đây là một tính từ dùng để chỉ các yếu tố tiền tệ.}

\vocabentry{Currency}
{/ˈkɜːrənsi/ (n)}
{A system of money in general use in a particular country.}
{Fluctuations in currency values can significantly impact international trade.}
{Đây là một danh từ dùng để chỉ đơn vị tiền tệ.}

\vocabentry{Exchange rate}
{/ɪksˈtʃeɪndʒ reɪt/ (n)}
{The value of one currency for the purpose of conversion to another.}
{A stable exchange rate is important for businesses involved in global export.}
{Đây là một cụm danh từ dùng để chỉ tỷ giá hối đoái.}

\vocabentry{Inflation}
{/ɪnˈfleɪʃn/ (n)}
{A general increase in prices and fall in the purchasing value of money.}
{High global energy prices have led to an increase in inflation in many countries.}
{Đây là một danh từ dùng để chỉ sự lạm phát.}

\vocabentry{Deflation}
{/ˌdiːˈfleɪʃn/ (n)}
{Reduction of the general level of prices in an economy.}
{Deflation can be just as damaging to the global economy as high inflation.}
{Đây là một danh từ dùng để chỉ sự xẹp phát/giảm phát.}

\vocabentry{Recession}
{/rɪˈseʃn/ (n)}
{A period of temporary economic decline during which trade and industrial activity are reduced.}
{The global recession of 2008 led to a sharp decrease in international trade.}
{Đây là một danh từ dùng để chỉ sự suy thoái kinh tế.}

\vocabentry{Depression}
{/dɪˈpreʃn/ (n)}
{A long and severe recession in an economy or market.}
{The Great Depression of the 1930s showed the dangers of an unstable global financial system.}
{Đây là một danh từ dùng để chỉ cuộc đại khủng hoảng kinh tế.}

\vocabentry{Growth rate}
{/ɡroʊθ reɪt/ (n)}
{The percentage increase in a country's real GDP.}
{Emerging economies often have much higher growth rates than developed ones.}
{Đây là một cụm danh từ dùng để chỉ tốc độ tăng trưởng.}

\vocabentry{GDP}
{/ˌdʒiː diː ˈpiː/ (n)}
{Gross Domestic Product; the total value of goods produced and services provided in a country during one year.}
{The United States has the largest GDP in the world.}
{Đây là một danh từ viết tắt dùng để chỉ tổng sản phẩm quốc nội.}

\vocabentry{GNI}
{/ˌdʒiː en ˈaɪ/ (n)}
{Gross National Income; the total domestic and foreign output claimed by residents of a country.}
{GNI per capita is a common measure of a country's standard of living.}
{Đây là một danh từ viết tắt dùng để chỉ tổng thu nhập quốc gia.}

\vocabentry{Wealth gap}
{/welθ ɡæp/ (n)}
{The unequal distribution of assets among residents of a country.}
{The global wealth gap between the richest 1\% and the rest of the world is widening.}
{Đây là một cụm danh từ dùng để chỉ khoảng cách giàu nghèo.}

\vocabentry{Poverty line}
{/ˈpɑːvərti laɪn/ (n)}
{The estimated minimum level of income needed to secure the necessities of life.}
{Millions of people in developing countries still live below the poverty line.}
{Đây là một cụm danh từ dùng để chỉ ngưỡng nghèo.}

\vocabentry{Humanitarian}
{/hjuːˌmænɪˈteriən/ (adj/n)}
{Concerned with or seeking to promote human welfare.}
{Globalization has made it easier to coordinate global humanitarian aid during disasters.}
{Đây là một tính từ/danh từ dùng để chỉ tính nhân đạo hoặc người làm nhân đạo.}

\vocabentry{Philanthropy}
{/fɪˈlænθrəpi/ (n)}
{The desire to promote the welfare of others, expressed by the generous donation of money to good causes.}
{Global philanthropy plays a key role in funding education and healthcare in poor nations.}
{Đây là một danh từ dùng để chỉ hoạt động từ thiện/nhân đạo.}

\vocabentry{Diversity}
{/daɪˈvɜːrsəti/ (n)}
{The state of being diverse; variety.}
{Cultural diversity is a source of strength and creativity for our world.}
{Trong văn hóa, đây là sự đa dạng.}

\vocabentry{Inclusion}
{/ɪnˈkluːʒn/ (n)}
{The action or state of including or of being included within a group or structure.}
{Global companies are focusing on diversity and inclusion in their hiring practices.}
{Đây là một danh từ dùng để chỉ sự hòa nhập/bao trùm.}

\vocabentry{Stereotype}
{/ˈsteriətaɪp/ (n)}
{A widely held but fixed and oversimplified image or idea of a particular type of person or thing.}
{Increased cultural exchange helps to break down global stereotypes.}
{Đây là một danh từ dùng để chỉ định kiến hoặc khuôn mẫu.}

\vocabentry{Prejudice}
{/ˈpredʒədɪs/ (n)}
{Preconceived opinion that is not based on reason or actual experience.}
{Education is the most powerful tool for fighting prejudice and intolerance.}
{Đây là một danh từ dùng để chỉ thành kiến.}

\vocabentry{Tolerance}
{/ˈtɑːlərəns/ (n)}
{The ability or willingness to tolerate something, in particular the existence of opinions or behavior that one does not necessarily agree with.}
{A globalized world requires a high level of religious and cultural tolerance.}
{Đây là một danh từ dùng để chỉ sự bao dung.}

\vocabentry{Discrimination}
{/dɪˌskrɪmɪˈneɪʃn/ (n)}
{The unjust or prejudicial treatment of different categories of people or things.}
{International laws forbid discrimination based on race, gender, or nationality.}
{Đây là một danh từ dùng để chỉ sự phân biệt đối xử.}

\vocabentry{Legacy}
{/ˈleɡəsi/ (n)}
{An amount of money or property left to someone in a will; also the long-lasting impact of events.}
{The legacy of colonialism continues to influence political stability in Africa.}
{Đây là một danh từ dùng để chỉ di sản hoặc hệ quả để lại.}

\vocabentry{Heritage site}
{/ˈherɪtɪdʒ saɪt/ (n)}
{A place that has been officially recognized for its cultural, historical, or natural importance.}
{UNESCO works to protect heritage sites all over the world.}
{Đây là một cụm danh từ dùng để chỉ di tích di sản thế giới.}

\vocabentry{Preservation}
{/ˌprezərˈveɪʃn/ (n)}
{The action of preserving something.}
{The preservation of endangered languages is a major global concern.}
{Đây là một danh từ dùng để chỉ sự bảo tồn văn hóa hoặc thiên nhiên.}

\vocabentry{Restoration}
{/ˌrestəˈreɪʃn/ (n)}
{The action of returning something to a former owner, place, or condition.}
{International teams worked together on the restoration of ancient monuments in Iraq.}
{Đây là một danh từ dùng để chỉ sự phục hồi.}

\vocabentry{Conservation}
{/ˌkɑːnsərˈveɪʃn/ (n)}
{The protection of plants and animals and natural areas.}
{Global conservation efforts are needed to save the world's oceans from pollution.}
{Đây là một danh từ dùng để chỉ sự bảo tồn tài nguyên thiên nhiên.}

\vocabentry{Awareness}
{/əˈwernəs/ (n)}
{Knowledge or perception of a situation or fact.}
{Social media has raised global awareness of human rights abuses.}
{Đây là một danh từ dùng để chỉ sự nhận thức.}

\vocabentry{Campaign}
{/kæmˈpeɪn/ (n)}
{An organized course of action to achieve a goal.}
{The WHO launched a global campaign to eliminate polio.}
{Đây là một danh từ dùng để chỉ chiến dịch toàn cầu.}

\vocabentry{Impact}
{/ˈɪmpækt/ (n/v)}
{The action of one object coming forcibly into contact with another; also a marked effect or influence.}
{The impact of globalization on local cultures is a subject of much debate.}
{Đây là một danh từ/động từ dùng để chỉ tác động hoặc ảnh hưởng.}

\vocabentry{Outcome}
{/ˈaʊtkʌm/ (n)}
{The way a thing turns out; a consequence.}
{The outcome of the trade talks will affect prices for consumers worldwide.}
{Đây là một danh từ dùng để chỉ kết quả hoặc hệ quả.}

\vocabentry{Consequence}
{/ˈkɑːnsəkwens/ (n)}
{A result or effect of an action or condition.}
{Job losses in manufacturing are a common consequence of offshoring.}
{Đây là một danh từ dùng để chỉ hậu quả hoặc hệ quả.}

\vocabentry{Trend}
{/trend/ (n)}
{A general direction in which something is developing or changing.}
{One of the major trends in globalization is the rise of the digital nomad.}
{Đây là một danh từ dùng để chỉ xu hướng toàn cầu.}

\vocabentry{Perspective}
{/pərˈspektɪv/ (n)}
{A particular attitude toward or way of regarding something; a point of view.}
{Traveling to different countries can provide a new perspective on your own culture.}
{Đây là một danh từ dùng để chỉ góc nhìn hoặc quan điểm cá nhân.}

\vocabentry{Viewpoint}
{/ˈvjuːpɔɪnt/ (n)}
{A person's opinion or point of view.}
{It is important to consider multiple viewpoints when analyzing global issues.}
{Đây là một danh từ đồng nghĩa với quan điểm.}

\vocabentry{Scenario}
{/səˈnærioʊ/ (n)}
{A postulated sequence or development of events.}
{Scientists have developed several scenarios for the future of the global climate.}
{Đây là một danh từ dùng để chỉ kịch bản có thể xảy ra trong tương lai.}

\vocabentry{Forecast}
{/ˈfɔːrkæst/ (n/v)}
{A prediction or estimate of future events.}
{Economic forecasts suggest a slowdown in global trade for the coming year.}
{Đây là một danh từ/động từ dùng để chỉ sự dự báo.}

\vocabentry{Prediction}
{/prɪˈdɪkʃn/ (n)}
{A thing predicted; a forecast.}
{Many early predictions about the "death of the nation-state" have proven incorrect.}
{Đây là một danh từ dùng để chỉ sự dự đoán.}

\vocabentry{Analysis}
{/əˈnæləsɪs/ (n)}
{Detailed examination of the elements or structure of something.}
{A thorough analysis of global markets is necessary before making an investment.}
{Đây là một danh từ dùng để chỉ sự phân tích.}

\vocabentry{Statistics}
{/stəˈtɪstɪks/ (n)}
{The practice or science of collecting and analyzing numerical data in large quantities.}
{Recent statistics show that global internet usage has passed 60\%.}
{Đây là một danh từ dùng để chỉ số liệu thống kê.}

\vocabentry{Evidence}
{/ˈevɪdəns/ (n)}
{The available body of facts or information indicating whether a belief or proposition is true or valid.}
{There is clear evidence that globalization has reduced extreme poverty worldwide.}
{Đây là một danh từ dùng để chỉ bằng chứng học thuật.}

\vocabentry{Correlation}
{/ˌkɔːrəˈleɪʃn/ (n)}
{A mutual relationship or connection between two or more things.}
{There is a strong correlation between economic openness and GDP growth.}
{Đây là một danh từ dùng để chỉ sự tương quan.}

\vocabentry{Significant}
{/sɪɡˈnɪfɪkənt/ (adj)}
{Sufficiently great or important to be worthy of attention; noteworthy.}
{Globalization has caused significant changes in the way we live and work.}
{Đây là một tính từ dùng để chỉ sự quan trọng hoặc đáng kể.}

\vocabentry{Essential}
{/ɪˈsenʃl/ (adj)}
{Absolutely necessary; extremely important.}
{International cooperation is essential for maintaining a stable global economy.}
{Đây là một tính từ dùng để chỉ sự thiết yếu.}

\vocabentry{Fundamental}
{/ˌfʌndəˈmentl/ (adj)}
{Forming a necessary base or core; of central importance.}
{The internet has caused a fundamental shift in global communications.}
{Đây là một tính từ dùng để chỉ tính chất nền tảng hoặc cơ bản.}

\vocabentry{Drastic}
{/ˈdræstɪk/ (adj)}
{Likely to have a strong or far-reaching effect; radical and extreme.}
{Some countries have taken drastic measures to protect their domestic industries.}
{Đây là một tính từ dùng để chỉ sự mạnh mẽ hoặc quyết liệt.}

\vocabentry{Gradual}
{/ˈɡrædʒuəl/ (adj)}
{Taking place or progressing slowly or by degrees.}
{The integration of the global economy has been a gradual process over several decades.}
{Đây là một tính từ dùng để chỉ sự dần dần, từ từ.}

\vocabentry{Rapid}
{/ˈræpɪd/ (adj)}
{Happening in a short time or at a great rate.}
{The rapid growth of emerging markets is changing the global balance of power.}
{Đây là một tính từ dùng để chỉ sự nhanh chóng.}

\vocabentry{Unprecedented}
{/ʌnˈpresɪdentɪd/ (adj)}
{Never done or known before.}
{We are living in an era of unprecedented global connectivity.}
{Đây là một tính từ dùng để chỉ tính chất chưa từng có tiền lệ.}

\vocabentry{Widespread}
{/ˌwaɪdˈspred/ (adj)}
{Found or distributed over a large area or number of people.}
{There is widespread concern about the impact of globalization on local jobs.}
{Đây là một tính từ dùng để chỉ sự lan rộng hoặc phổ biến.}

\vocabentry{Controversial}
{/ˌkɑːntrəˈvɜːrʃl/ (adj)}
{Giving rise or likely to give rise to public disagreement.}
{The benefits of economic globalization remain a highly controversial topic.}
{Đây là một tính từ dùng để chỉ tính tranh cãi.}

\vocabentry{Paradigm}
{/ˈpærədaɪm/ (n)}
{A typical example or pattern of something; a model.}
{Globalization has created a new paradigm for how international business is conducted.}
{Đây là một danh từ dùng để chỉ một mô hình hoặc hình mẫu chuẩn mực.}

